\documentclass[11pt]{article}
\usepackage[margin=1in]{geometry}
\usepackage{amsmath,amssymb}
\usepackage[round]{natbib}
\usepackage{authblk}
\usepackage{hyperref}
\title{Robust Change-Point Detection for Heavy-Tailed Sensor Streams}
\author[1]{Yuzuki Tachibana}
\author[2]{Nora Lindqvist}
\author[1]{Daichi Morioka}
\affil[1]{Department of Statistics, Kitamachi University, Japan}
\affil[2]{Institute of Fictional Statistics, University of Nowhere, Freedonia}
\date{\today}
\begin{document}
\maketitle
\begin{abstract}
Change-point detection methods often assume light-tailed noise, but measurements from low-cost sensors frequently contain outliers and heavy tails.
We propose a fictional robust procedure that combines a median-of-means estimator with a CUSUM-type statistic and adapts its threshold to the unknown tail index.
We prove that the procedure controls the false alarm rate under only a second-moment assumption and that its detection delay is optimal up to logarithmic factors.
Simulations show that it outperforms existing robust methods when the noise has infinite fourth moment, while losing little efficiency under Gaussian noise.
An application to a fictional air-quality network illustrates its practical value.
\end{abstract}
\noindent\textbf{Keywords:} change-point detection; robust statistics; heavy tails; median of means
\section{Introduction}
Detecting abrupt changes in a data stream is a classical problem with applications from quality control to environmental monitoring \citep{tachibana2021,lindqvist2022}.
Most procedures rely on sums of observations and therefore inherit their sensitivity to outliers.
Robust alternatives exist, but they typically require knowledge of the noise distribution or a tuning parameter that is hard to choose \citep{morioka2023}.

In this paper we propose a procedure that is robust to heavy tails and does not require such knowledge.
The main idea is to replace sample means by median-of-means estimators computed on blocks whose size grows with the length of the monitoring window.

\section{Method}
Let $X_1, X_2, \dots$ be independent observations.
\bibliographystyle{plainnat}
\bibliography{refs}
\end{document}
